\section{Useful-Compute Markets}
\label{sec:markets}

Hanzo Network exposes inference, training, evaluation, and capacity markets
through one A-Chain job format. A market determines who requests the work and
who pays for it. PoAI consensus determines whether the bound execution was
actually completed and whether a mining reward may be finalized. These are
separate decisions.

\subsection{Canonical mining job}

Before execution, A-Chain commits a job containing

\begin{equation}
\begin{split}
J=H(&\text{job id},\text{requester},\text{consumer},\text{model root},
\text{input root},\\
&\text{execution policy},\text{proof policy},\text{reward budget},
\text{expiry},\text{destination},\text{recipient}).
\end{split}
\end{equation}

The job must be admitted before a miner commits its result. This prevents the
miner from choosing an arbitrary completed computation and declaring it useful
after the fact. ``Useful'' means useful under an explicit demand or governed
network-service policy; it does not claim universal truth or social value.

\subsection{Execution and commitment}

A CPU or GPU provider runs the specified integer program and commits its output,
typed graph transcript, metering, provider identity, and runtime identifier. The
canonical path defines tensor layout, serialization, accumulator width,
overflow, rounding, saturation, routing, and tie-breaking. An honest CPU and GPU
must therefore produce identical committed integer values.

The commitment precedes any random challenge. A job-specific nonce, policy
epoch, and execution commitment form the global work nullifier, so the same work
authorization cannot be replayed into a second chain or reward.

\subsection{CPU validation}

PoAI does not require every validator to repeat a GPU-scale workload. The active
proof policy chooses one of four paths:

\begin{itemize}[leftmargin=1.2em]
  \item \textbf{Exact replay}: appropriate for small jobs; an ordinary CPU
        executes the canonical integer program.
  \item \textbf{Optimistic integer proof}: the provider publishes a Merkle-rooted
        transcript; watchers re-execute and may challenge any invalid operation.
        For an opened matrix product $C=AB$, a CPU checks
        $A(Br)=Cr$ over $\mathbb{F}_{2^{61}-1}$ after $r$ is derived from a
        post-commitment beacon. The check is quadratic rather than cubic.
  \item \textbf{Succinct proof}: an audit- and activation-gated P3Q circuit may
        replace the challenge game. A generic STARK/FRI library is not, by
        itself, a production PoAI verifier.
  \item \textbf{Confidential execution}: an approved TEE profile may add vendor
        provenance for private prompts, weights, or outputs. It is an optional
        trust policy and not the definition of PoAI.
\end{itemize}

Freivalds supplies local matrix-product soundness, not automatic whole-model
soundness. The complete policy must also enforce graph wiring, challenge
coverage, transcript availability, metering integrity, and a finality delay that
prevents minting while a challenge remains open.

\subsection{Inference market}

A requester posts a model, input commitment, output contract, latency bound, and
maximum demand fee. Providers bid or accept the open price. On successful PoAI
finality, the consumer receives the output, the provider receives its demand
fee, and the receipt may receive a mining reward under the current era budget.
The mining reward is not a substitute for the customer's payment; it rewards
making the execution independently accountable.

\subsection{Training and evaluation markets}

A training job additionally commits the parent model, dataset root, optimizer
state, loss, deterministic randomness, checkpoint cadence, and target metric.
Only training graphs expressible in the activated deterministic operation set
are proof-bearing. Floating-point distributed training is not silently treated
as consensus deterministic.

An evaluation job commits the model, dataset or sealed test root, metric, and
aggregation policy. A subjective judge quorum may earn a demand fee for a signed
assessment, but a quorum signature alone cannot authorize mining subsidy.

\subsection{Capacity market}

The capacity market routes open jobs to advertised CPU/GPU inventory. Hardware
class, memory, region, price, and availability affect matching, but raw rented
seconds do not mint AI. A provider earns mining rewards only for an admitted job
whose proof and metering reach PoAI finality.

\subsection{Existing cloud-provider adapter}

The same job format is intended for hyperscalers and independent operators. A
cloud adapter captures the model/input commitment and execution policy when the
customer submits an ordinary workload, records the integer transcript alongside
normal execution, and sends the proof commitment to A-Chain. The provider keeps
its existing invoice revenue and may receive AI when PoAI consensus finalizes
the work. Customer opt-in, confidentiality policy, and export controls remain
part of job admission.

This gives existing clouds a direct reason to participate: new revenue with no
need to abandon their scheduler, a portable cryptographic service receipt for
customers, and access to decentralized overflow capacity. The network gains
real workload volume instead of synthetic benchmark puzzles.

\subsection{Anti-farming rules}

The reward policy must reject circular or unverifiable demand. At minimum:

\begin{itemize}[leftmargin=1.2em]
  \item related requester/provider control is declared and may be capped or
        excluded;
  \item demand fees and network-service budgets are bounded and non-recyclable
        under the active policy;
  \item metering is committed and objectively challengeable;
  \item output copying without the bound transcript fails the challenge;
  \item every finalized receipt consumes one global A-Chain nullifier;
  \item aggregate payouts never exceed the current halving-era allowance.
\end{itemize}

\subsection{Spot and reserved service}

Providers may quote spot or reserved prices. Spot jobs trade cancellation risk
for lower price; reserved jobs escrow a forward capacity commitment. Both use
the same A-Chain job, proof, nullifier, receipt, and settlement path. Pricing is
a market concern; proof validity is not.
